\section{Vérification fonctionnelle du firmware}
\label{sec:res_firmware}

Cette phase de validation parcourt le firmware via son interface utilisateur sur la carte finale, écran par écran et fonction par fonction. Elle se distingue de la validation technique de la section~\ref{sec:soft_validation} qui testait les mécanismes internes par injection de fautes durant le développement. Ici, aucun instrument n'est utilisé et chaque point est jugé sur le comportement observable de l'appareil. Les 69 points de contrôle et leurs résultats figurent à la section~\ref{ann:mesures_sw} de l'annexe.

Le périmètre de test couvre le menu principal, les quatre pages de statistiques, le navigateur de stockage, l'écran de lecture, les réglages, divers cas dégradés, la robustesse au redémarrage et l'absence de bruits parasites. Sur les 69 vérifications, 66 sont conformes. Les trois écarts observés n'affectent pas la lecture audio en fonctionnement standard.

Deux non-conformités concernent la gestion d'erreur à l'ouverture de fichier. Un fichier WAV tronqué ou vide, ainsi qu'un fichier d'extension supportée mais de format incompatible, lancent la lecture au lieu d'afficher un message d'erreur. Le pipeline audio restant fonctionnel et l'appareil ne se bloquant pas, la conséquence se limite à l'absence de retour utilisateur. Bien que la reconnaissance de format décrite à la section~\ref{sec:transport} identifie correctement le conteneur par sa signature, aucune vérification ultérieure ne valide la cohérence du contenu avec cette signature.

La troisième non-conformité porte sur un bruit de commutation résiduel. Aucune impulsion n'est audible à l'allumage, à l'extinction, au début ou à la fin d'une lecture, ni lors du changement de sortie. Un léger glitch subsiste cependant lors du changement de piste. Son origine n'a pas été formellement établie faute de temps. L'hypothèse la plus probable est logicielle car les quatre transitions propres exploitent la séquence de sourdine du convertisseur et de mise en veille de l'amplificateur décrite à la section~\ref{sec:schema}. Le chemin de reconfiguration emprunté lors du changement de piste ne suit pas cette séquence.